
\lesson{22/10}
The momentum operator can be derived in a similar fashion to the Hamiltonian:
$$\hat P^i = \int \dd^3x \hat T\indices{^i_0} = \int \dd ^3x\hat\pi\pd^i\hat\phi$$
substituting in the F.Ts of the two fields and using the commutation relations;
using the normal ordering this yields
$$:\hat P^i: = \int \dd^3k k^i\hat a^\dagger_k \hat a_k$$
which is similar in form to $\hat H$, since $\omega_k=k^0$.
\subsection{Fock Space}
The space of the states on which the operators are defined is 
a "reordered" Hilbert Space:
$$\mathcal{H}_\text{Fock} = \mathcal{H}_1\oplus\mathcal{H}_2\oplus...$$
where $\mathcal{H}_1$ is the Hilbert subspace in which the basis is made
by all the stated with $1$ particle and
$$\mathcal{H}_i = \underbrace{\mathcal{H}_1\otimes\mathcal{H}_1\otimes...}_
{i \text{ times}}.$$
The $\mathcal{H}_0$ space is trivial, as it contains only the vacuum state
$\ket{0}$.\\

The creation and annihilation operators add or remove particles from the state:
\begin{gather*}
\hat a(\mathbf{k})\ket{0}=0\\
\hat a^\dagger(\mathbf{k})\ket{0} := \ket{\mathbf{k}}\in\mathcal{H}_1\\
\ket{\mathbf{k_1,k_2,...k_n}} = \hat a^\dagger(\mathbf{k_1});
\hat a^\dagger(\mathbf{k_2})...\hat a^\dagger(\mathbf{k_n})\ket{0},
\end{gather*}
which represents a state with $n$ particles of momentum $\mathbf{k_i}$: if these
momenta are equal, the shorthand notation is
$$\ket{n(\mathbf{k})}=\frac{(\hat a^\dagger(\mathbf{k}))}{\sqrt{n!}}\ket{0}.$$
The number operator is defined as
$$\hat n(\mathbf{k}) = \hat a^\dagger(\mathbf{k})\hat a(\mathbf{k})$$
and its eigenstates are $\ket{n(\mathbf{k})}$:
$$\hat n(\mathbf{k})\ket{n(\mathbf{k})}=n\ket{n(\mathbf{k})}.$$

The orthogonality condition
$$\bra{\mathbf{k}}\ket{\mathbf{k'}}=\delta^3(\mathbf{k-k'})$$
is not Lorentz invariant, so the normalization is set to be
$$\int \dd^3k\frac{1}{(2\pi)^32\omega_k}\ket{k}\bra{k}=\mathbb{I}$$
where $\ket{k}$ is the state with 4-momentum $k$ and differs
from $\ket{\mathbf{k}}$ only by a normalization
$$\ket{k} = (2\pi)^{\frac{3}{2}}\sqrt{2\omega_k}\ket{\mathbf{k}}.$$

How does the field operator act on the vacuum?
$$\hat\phi\ket{0} = \int \dd^3k\mathcal{N}\left[\hat a_k\ket{0}e^{-ikx}
    +\hat a_k^\dagger\ket{0}e^{ikx}
\right] = \int \dd^3k \mathcal{N} e^{ikx}\ket{\mathbf{k}}$$
this can be interpreted as the field exciting the vacuum creating
a particle for each and every oscillating mode (given by $\mathbf{k}$)

Attaching a $\bra{k'}$ on the left yields a transition
amplitude between $\mathbf{k}$ and $\mathbf{k'}$ 
$$\bra{k'}\hat\phi(x^\mu)\ket{0} = \int \dd^3k \frac{\sqrt{2\omega_{k'}}}
{\sqrt{2\omega_k}}e^{ikx}\delta^3(\mathbf{k-k'}) = e^{ik'x}$$

\textbf{Note} With this formalization, the Bose-Einstein Statistic for scalar
particles comes for free: consider the creation of two particles with different
momenta
$$a^\dagger(k')a^\dagger(k)\ket{0} = a^\dagger(k)a^\dagger(k')\ket{0}\neq0,$$
since $\left[a^\dagger(k),a^\dagger(k')\right] = 0$. The exchange in the
order of creation of the particles is irrelevant (this will not be the case
for the Dirac Field). Moreover the commutation relation does not forbid the
successive creation of $n$ particles with the same momentum $k^i$:
$$\forall n,\; (a^\dagger(k))^n\ket{0}\neq0$$
effectively resembling the Bose-Einstein statistic.
\subsection{Canonical Quantization of the charged scalar field}
The Lagrangian of the charged scalar field is
$$\mathcal{L}=\pd_\mu\phi^\dagger\pd^\mu\phi - m^2\phi\phi^\dagger;$$
the conjugated fields are two:
\begin{gather*}
    \pi_\phi= \pdv[]{\mathcal{L}}{\pd_0\phi} = {\dot \phi}^\dagger\\
    \pi_{\phi^\dagger}= \pdv{\mathcal{L}}{\pd_0\phi^\dagger} ={\dot \phi}
\end{gather*}
The Hamiltonian density is defined as
$$\mathcal{H}=\pi_\phi\dot\phi + \pi_{\phi^\dagger}\dot\phi^\dagger -\mathcal{L}=
\dot\phi^\dagger\dot\phi - \pd_i\phi^\dagger\pd^i\phi + m^2\phi^\dagger\phi:$$
now the fields are promoted to operators, with commutation relations
$$\left[\phi(t,x),\pi_\phi(t,y)\right]=\left[\phi^\dagger(t,x),\pi_{\phi^\dagger}
\right]=i\delta^3(\mathbf{x-y})$$
and all the other are zero.

The Fourier transforms of the two fields are
\begin{gather*}
    \phi(x^\mu) = \int \frac{\dd^3k}{(2\pi)^{\frac{3}{2}}\sqrt{2\omega_k}}
    \left[a_ke^{-ikx}+b^\dagger_k e^{ikx}\right]\\
    \phi^\dagger(x^\mu) = \int \frac{\dd^3k}{(2\pi)^{\frac{3}{2}}\sqrt{2\omega_k}}
    \left[a^\dagger_ke^{ikx}+b_k e^{-ikx}\right]\\
\end{gather*}
\textbf{Note} There are two sets of operators because the charged scalar field
is composed of two different real scalar fields.

And so, $a$ are associated to particles, while $b$ to antiparticles.

The canonical commutation relations yield
$$\left[a_k,a^\dagger_{k'}\right] = \left[b_k,b^\dagger_{k'}\right]=\delta^3(
\mathbf{k-k'})$$
and all others are zero: this is the algebraic structure of two decoupled
harmonic oscillators.

With a big ugly calculation, the normal ordered Hamiltonian turns out to be
$$:H: = \int \dd ^3k\omega_k\left(a_k^\dagger a_k+b^\dagger_k b_k\right)=
\int \dd ^3x \omega_k\left(n^a(\mathbf{k})+n^b(\mathbf{k})\right).$$
\subsubsection{Noether currents and charges}
Recall that the charged scalar field Lagrangian is invariant under the
transformation
$$\phi\to\phi'=e^{i\alpha}\phi:$$
the Noether current associated to this symmetry is an operator and it is
$$J^\mu = \phi^\dagger\pd_\mu\phi - \pd_\mu\phi^\dagger \phi$$
and the conserved charge is
$$Q = \int \dd^3x J^0 = \int \dd^3x (\phi^\dagger\dot\phi-\dot\phi^\dagger\phi)$$
replacing the fields' F.Ts in this expression yields
$$Q=\int \dd^3k \left(a^\dagger_ka_k - b^\dagger_kb_k\right) = 
\int \dd^3k \left(n^a(\mathbf{k})-n^b(\mathbf{k})\right).$$
This can be interpreted as the \textit{total charge} of the system, with the "a"
particles have positive charge, while the "b" particles have negative charge:
the conserved quantity is then the total (Abelian) charge and this conservation
means that creating an "a" particle implies creating a "b" particle.
\subsection{Canonical Quantization of the Dirac Field}
The Lagrangian for the spinorial field is
$$\mathcal{L}=\bar\psi(i\gamma^\mu\pd_\mu-m)\psi$$
The conjugated fields are
\begin{gather*}
\pi_\psi = \pdv[]{\mathcal{L}}{\dot\psi}=i\psi^\dagger,\qquad
\pi_{\bar\psi} = \pdv[]{\mathcal{L}}{\dot{\bar\psi}}=0
\end{gather*}
and so the Hamiltonian density is
$$\mathcal{H} = \pi_\psi\dot\psi + \pi_{\bar\psi}\dot{\bar\psi} -\mathcal{L}=
\bar\psi(-i\gamma^i\pd_i+m)\psi$$
using the Dirac equation to express $i\gamma^i\pd_i = m-i\gamma^0\pd_0$ yields
$$\mathcal{H} = i\psi^\dagger\dot\psi.$$
Now the commutation relations need to be set between $\psi,\bar\psi$ and their
conjugated fields: since these are fermionic fields, anticommutators are expected
(since a theory with commutators does not conserve energy):
$$
\left\{\psi_s(t,x),\pi_{\psi,r}(t,y)\right\}=i\delta^3(\mathbf{x-y})\delta_{rs}\to
\left\{\psi_s(t,x),\psi_r^\dagger(t,y)\right\}=\delta^3(\mathbf{x-y})\delta_{rs},
\;\text{with }r,s=1,2.
$$
\textbf{Note} Since the fields are spinors, these anticommutation relations also
depend on the spin of the field.
